% Options for packages loaded elsewhere
\PassOptionsToPackage{unicode}{hyperref}
\PassOptionsToPackage{hyphens}{url}
\documentclass[
]{article}
\usepackage{xcolor}
\usepackage{amsmath,amssymb}
\setcounter{secnumdepth}{-\maxdimen} % remove section numbering
\usepackage{iftex}
\ifPDFTeX
  \usepackage[T1]{fontenc}
  \usepackage[utf8]{inputenc}
  \usepackage{textcomp} % provide euro and other symbols
\else % if luatex or xetex
  \usepackage{unicode-math} % this also loads fontspec
  \defaultfontfeatures{Scale=MatchLowercase}
  \defaultfontfeatures[\rmfamily]{Ligatures=TeX,Scale=1}
\fi
\usepackage{lmodern}
\ifPDFTeX\else
  % xetex/luatex font selection
\fi
% Use upquote if available, for straight quotes in verbatim environments
\IfFileExists{upquote.sty}{\usepackage{upquote}}{}
\IfFileExists{microtype.sty}{% use microtype if available
  \usepackage[]{microtype}
  \UseMicrotypeSet[protrusion]{basicmath} % disable protrusion for tt fonts
}{}
\makeatletter
\@ifundefined{KOMAClassName}{% if non-KOMA class
  \IfFileExists{parskip.sty}{%
    \usepackage{parskip}
  }{% else
    \setlength{\parindent}{0pt}
    \setlength{\parskip}{6pt plus 2pt minus 1pt}}
}{% if KOMA class
  \KOMAoptions{parskip=half}}
\makeatother
\usepackage{longtable,booktabs,array}
\newcounter{none} % for unnumbered tables
\usepackage{calc} % for calculating minipage widths
% Correct order of tables after \paragraph or \subparagraph
\usepackage{etoolbox}
\makeatletter
\patchcmd\longtable{\par}{\if@noskipsec\mbox{}\fi\par}{}{}
\makeatother
% Allow footnotes in longtable head/foot
\IfFileExists{footnotehyper.sty}{\usepackage{footnotehyper}}{\usepackage{footnote}}
\makesavenoteenv{longtable}
\setlength{\emergencystretch}{3em} % prevent overfull lines
\providecommand{\tightlist}{%
  \setlength{\itemsep}{0pt}\setlength{\parskip}{0pt}}
\usepackage{bookmark}
\IfFileExists{xurl.sty}{\usepackage{xurl}}{} % add URL line breaks if available
\urlstyle{same}
\hypersetup{
  hidelinks,
  pdfcreator={LaTeX via pandoc}}

\author{}
\date{}

\begin{document}

\section{CROWDLAB: Supervised learning to infer consensus labels and
quality scores for data with multiple
annotators}\label{crowdlab-supervised-learning-to-infer-consensus-labels-and-quality-scores-for-data-with-multiple-annotators}

Hui Wen Goh \(^{1}\) Ulyana Tkachenko \(^{1}\) Jonas Mueller \(^{1}\)

\section{Abstract}\label{abstract}

Real-world data for classification is often labeled by multiple
annotators. For analyzing such data, we introduce CROWDLAB, a
straightforward approach to utilize any trained classifier to estimate:
(1) A consensus label for each example that aggregates the available
annotations; (2) A confidence score for how likely each consensus label
is correct; (3) A rating for each annotator quantifying the overall
correctness of their labels. Existing algorithms to estimate related
quantities in crowdsourcing often rely on sophisticated generative
models with iterative inference. CROWDLAB instead uses a straightforward
weighted ensemble. Existing algorithms often rely solely on annotator
statistics, ignoring the features of the examples from which the
annotations derive. CROWDLAB utilizes any classifier model trained on
these features, and can thus better generalize between examples with
similar features. On real-world multi-annotator image data, our proposed
method provides superior estimates for (1)-(3) than existing algorithms
like Dawid-Skene/GLAD.

\section{1. Introduction}\label{introduction}

Training data for multiclass classification are often labeled by
multiple annotators, with some redundancy between annotators to ensure
high-quality labels. Such settings have been studied in crowdsourcing
research (Monarch, 2021b; Paun et al., 2018). There it is often assumed
that many annotators have labeled each example (Carpenter, 2008; Khetan
et al., 2018), but this can be prohibitively expensive. This paper
considers general settings where each example in the dataset is merely
labeled by at least one annotator, and each annotator labels many
examples (but still only a subset of the dataset). Each annotation
corresponds to the selection of one class \(y \in \{1, \ldots, K\}\)
which the annotator believes to be most appropriate for this example.

Certain classification models can be trained in a special manner to
account for the multiple labels per example (Nguyen et al., 2014;
Peterson et al., 2019), but this is rarely done in practical
applications. A common approach is to aggregate the labels for each
example into a single consensus label, e.g.~via majority-vote or
statistical crowdsourcing algorithms (Dawid and Skene, 1979). Any
classifier can then be trained on these consensus labels via
off-the-shelf code.

Here we propose a method \(^{1}\) that leverages any already-trained
classifier to: (1) establish accurate consensus labels, (2) estimate
their quality, and (3) estimate the quality of each annotator (Monarch,
2021c). The latter two aims help us determine which data is least
trustworthy and should perhaps be verified via additional annotation
(Bernhardt et al., 2022). CROWDLAB (Classifier Refinement Of
croWD-sourced LABels) is based on a straightforward weighted ensemble of
the classifier predictions and individual annotations. Weights are
assigned according to the (estimated) trustworthiness of each annotator
relative to the trained classifier. CROWDLAB is easy to
implement/understand, computationally efficient (non-iterative), and
extremely flexible. It works with any classifier and training procedure,
as well as any classification dataset (including those containing
examples only labeled by one annotator).

Motivations. Illustrating how many real-world multi-annotator datasets
look, Figure 1 shows a disparity in annotator quality as well as many
examples whose consensus label will be incorrect if we rely on majority
vote (nonetheless often done in practice due to its straightforward
appeal). Unsurprisingly, consensus labels are more likely to be
incorrect for those examples with fewer annotations. An effective method
to estimate consensus label quality should properly account for the
number of annotations an example has received, as well as the quality of
the annotators who selected these labels. Many of the examples whose
consensus label is wrong merely have a single annotation, which provides
little information. Using a trained classifier can help us better
generalize to such examples to estimate their labels' quality
(especially if the data contain other examples with similar feature
values). When incorporating a classifier, we

\textit{[Image omitted: images/09f254b5eb907d4145fe4b136b38fdd4f88aef51224660aa2ab675152b7ca208.jpg]}

boxplot

{\def\LTcaptype{none} % do not increment counter
\begin{longtable}[]{@{}ll@{}}
\toprule\noalign{}
Annotator Accuracy & Value \\
\midrule\noalign{}
\endhead
\bottomrule\noalign{}
\endlastfoot
Min & 0.1 \\
Q1 & 0.2 \\
Median & 0.9 \\
Q3 & 0.85 \\
Max & 1.0 \\
\end{longtable}
}

\textit{[Image omitted: images/b98b48d344addb196d72336f6a587e709edd06838d34f808ad0ce4bc129526c2.jpg]}

boxplot

{\def\LTcaptype{none} % do not increment counter
\begin{longtable}[]{@{}ll@{}}
\toprule\noalign{}
Consensus label accuracy & Number of annotations for example \\
\midrule\noalign{}
\endhead
\bottomrule\noalign{}
\endlastfoot
Correct & 10 \\
Correct & 9 \\
Correct & 8 \\
Correct & 7 \\
Correct & 6 \\
Correct & 5 \\
Correct & 4 \\
Correct & 3 \\
Correct & 2 \\
Correct & 1 \\
Correct & 0 \\
Wrong & 7 \\
Wrong & 6 \\
Wrong & 5 \\
Wrong & 4 \\
Wrong & 3 \\
Wrong & 2 \\
Wrong & 1 \\
Wrong & 0 \\
\end{longtable}
}

Figure 1. Statistics of our Hardest dataset, which has images annotated
by many actual humans. (a) Distribution over annotators showing the
overall accuracy of each annotator's chosen labels. (b) Distribution
over examples showing the number of annotations per example, grouped by
whether the majority-vote consensus label is correct or not. Accuracy is
measured against underlying ground-truth labels.

also wish to account for the accuracy and confidence of its estimates.
CROWDLAB offers a straightforward way to appropriately account for all
of these factors.

\section{2. Methods}\label{methods}

Consider a dataset sampled from (feature, class label) pairs \((X,Y)\)
that is comprised of: n examples, K classes, and m annotators in total.
We first establish some notation to formally describe our setting:
\(|J|\) denotes the cardinality of set \(J,1(\cdot)\) is an indicator
function emitting 1 if its condition is True and 0 otherwise, and
\([n]=\{1,2,\ldots,n\}\) indexes examples in the dataset. \(X_{i}\)
denotes the features of the ith example, which belongs to one class
\(Y_{i}\in[K]\) . This true class is unknown to us. For \(j\in[m]\) :
\(A_{j}\) denotes the jth annotator, and \(Y_{ij}\in[K]\) is the class
this annotator chose for \(X_{i}\) . \(Y_{ij}=\varnothing\) if \(A_{j}\)
did not label this particular example. Each example receives at most m
annotations, with many examples receiving fewer. \(\hat{Y}_{i}\) denotes
the consensus label for example i, representing our best estimate of its
true class \(Y_{i}\) . \(I_{j}:=\{i\in[n]:Y_{ij}\neq\varnothing\}\)
denotes the subset of examples labeled by \(A_{j}\) . We assume each
annotator has labeled multiple examples, i.e.~\(|I_{j}|>1\) .
\(J_{i}:=\{j\in[m]:Y_{ij}\neq\varnothing\}\) denotes the subset of
annotators that labeled \(X_{i}\) . Some examples may only be labeled by
a single annotator.

We assume some classifier model \(\mathcal{M}\) has been trained to
predict the given labels based on feature values. CROWD-LAB can be used
with any type of classifier \(\mathcal{M}\) (and training procedure), as
long as it outputs predicted class probabilities
\(\widehat{p}_{\mathcal{M}}(Y \mid X) \in \mathbb{R}^K\) estimating the
likelihood that example \(X\) belongs to each class \(k\) . To avoid
overfit predictions, we fit \(\mathcal{M}\) via cross-validation. This
provides held-out predictions
\(\widehat{p}_{\mathcal{M}}(Y_i \mid X_i)\) for each example in the
dataset (from a copy of \(\mathcal{M}\) which never saw \(X_i\) during
training). In our experiments, we simply train \(\mathcal{M}\) on
consensus labels derived via majority vote. But one could train the
classifier on any other set of consensus labels or even on the
individual labels from each annotator (simply duplicating
multiply-annotated examples in the training set). All methods considered
here that use \(\mathcal{M}\) will benefit from improvements in the
classifier's predictive accuracy. However CROWDLAB is the only method
that explicitly accounts for shortcomings of the classifier's
predictions (inevitable due to estimation error).

\section{2.1. Scoring Consensus
Quality}\label{scoring-consensus-quality}

We first outline methods to estimate our confidence that a given
consensus label for each example is correct. These quality estimates
\(q_{i} \in [0, 1]\) may be applied to any given label no matter which
method was used to establish consensus. Once we can estimate the quality
of any one label for each example, we estimate the best consensus label
under each method as the class with the highest consensus quality score.
This class can be identified efficiently for CROWDLAB. CROWDLAB combines
the complementary strengths of two basic estimators that we discuss
first.

Agreement (Monarch, 2021b). The fraction of annotators who agree with
consensus label (does not use a classifier).

\[
q _ {i} = \frac {1}{| \mathcal {J} _ {i} |} \sum_ {j \in \mathcal {J} _ {i}} \mathbb {1} \left(Y _ {i j} = \widehat {Y} _ {i}\right) \tag {1}
\]

Final consensus labels can be established via majority vote.

Label Quality Score (Kuan and Mueller, 2022). Instead of relying on the
annotators, one can rely on the classifier model via methods used to
evaluate labels in standard (singly-labeled) classification datasets.
This approach ignores information from individual annotators when
computing consensus quality:
\(q_{i} = L(\widehat{Y}_{i}, \widehat{p}_{\mathcal{M}}(Y_{i} \mid X_{i}))\)
.

Here \(L(Y, \hat{p}) \in [0, 1]\) is a label quality score which
quantifies our confidence that a particular label \(Y \in [K]\) is
correct for example X, given model-prediction \(\hat{p} \in R^{K}\)
estimating the likelihood that X belongs to each class. Our work uses
self-confidence as the label quality score:
\(L(Y, p) = \hat{p}(Y \mid X)\) . This simply represents the
model-estimated probability that the example belongs to its labeled
class. Kuan and Mueller (2022); Northcutt et al.~(2021b) found this to
be effective for scoring label errors in singly-labeled data based on
classifier predictions.

CROWDLAB (Classifier Refinement Of croWDsourced LABels). The
aforementioned approaches fail to consider both annotators and
classifier. Treating these as different predictors of an example's true
label, we take inspiration from prediction competitions where weighted
ensembling of predictors produces accurate and calibrated predictions.
CROWDLAB also employs the same label quality score for each consensus
label, but applies it to a different class probability vector which
modifies the prediction output by our classifier to account for the
individual annotations for an example:
\(q_i = L(\hat{Y}_i, \hat{p}_{\text{CR}}(Y_i \mid X_i, \{Y_{ij}\}))\) .

We estimate these class probabilities by means of a weighted ensemble
aggregation (Fakoor et al., 2021):

\[
\widehat {p} _ {\mathrm{CR}} (Y _ {i} \mid X _ {i}, \{Y _ {i j} \}) =
\]

\[
\frac {w _ {\mathcal {M}} \cdot \widehat {p} _ {\mathcal {M}} (Y _ {i} \mid X _ {i}) + \sum_ {j \in \mathcal {J} _ {i}} w _ {j} \cdot \widehat {p} _ {\mathcal {A} _ {j}} (Y _ {i} \mid \{Y _ {i j} \})}{w _ {\mathcal {M}} + \sum_ {j \in \mathcal {J} _ {i}} w _ {j}}
\]

Here \(\widehat{p}_{M} \in R^{K}\) is the probability of each class
predicted by our classifier, \(\widehat{p}_{A_{j}} \in R^{K}\) is a
similar likelihood vector treating each annotator's label as a
probabilistic ``prediction'', and \(w_{j}, w_{M} \in R\) are weights to
account for the (estimated) relative trustworthiness of each annotator
and our classifier. Our estimation procedure for these weights ensures
\(w_{M}\) is smaller if our classifier was poorly trained and \(w_{j}\)
is smaller for the annotators who give less accurate labels overall.

To present the remaining details, we first define a likelihood parameter
P as the average annotator agreement, across examples with more than one
annotation. P estimates the probability that an arbitrary annotator's
label will match the majority-vote consensus label for an arbitrary
example.

\[
P = \frac {1}{| \mathcal {I} _ {+} |} \sum_ {i \in \mathcal {I} _ {+}} \frac {1}{| \mathcal {J} _ {i} |} \sum_ {j \in \mathcal {J} _ {i}} \mathbb {1} (Y _ {i j} = \hat {Y} _ {i})
\]

where \(I_{+} := \{i \in [n] : |J_{i}| > 1\}\) (2)

We then simply define our per annotator predicted class probability
vector used in (2.1) to be:

\[
\widehat {p} _ {\mathcal {A} _ {j}} \left(Y _ {i} = k \mid \left\{Y _ {i j} \right\}\right) = \left\{ \begin{array}{l l} P & \text { when } Y _ {i j} = k \\ \frac {1 - P}{K - 1} & \text { when } Y _ {i j} \neq k \end{array} \right. \tag {3}
\]

This likelihood is shared across annotators and only involves a single
parameter P, easily estimated from limited data. P is a simple estimate
of the accuracy of labels from a typical annotator. Note that including
singly-annotated examples in (2) would bias P. This likelihood
facilitates comparing classifier outputs against outputs from the
typical annotator.

Now we detail how to estimate the trustworthiness weights
\(w_{j}, w_{M}\) . Let \(s_{j}\) represent annotator j's agreement with
other annotators who labeled the same examples:

\[
s _ {j} = \frac {\sum_ {i \in \mathcal {I} _ {j}} \sum_ {\ell \in \mathcal {J} _ {i} , \ell \neq j} \mathbb {1} \left(Y _ {i j} = Y _ {i \ell}\right)}{\sum_ {i \in \mathcal {I} _ {j}} \left(\left| \mathcal {J} _ {i} \right| - 1\right)} \tag {4}
\]

Let \(A_{\mathcal{M}}\) be the (empirical) accuracy of our classifier
with respect to the majority-vote consensus labels over the examples
with more than one annotation:

\[
A _ {\mathcal {M}} = \frac {1}{| \mathcal {I} _ {+} |} \sum_ {i \in \mathcal {I} _ {+}} \mathbb {1} \left(Y _ {i, \mathcal {M}} = \hat {Y} _ {i}\right) \tag {5}
\]

Here
\(Y_{i,\mathcal{M}} := \arg \max_{k} \widehat{p}_{\mathcal{M}}(Y_{i} = k \mid X_{i}) \in [K]\)
is the class predicted by our model for \(X_{i}\) . \(A_{M}\) and
\(s_{j}\) from (4) are analogous accuracy estimates for our classifier
and individual annotators. Both are computed with only the
multiply-annotated examples, since majority-vote consensus labels are
more reliable for this subset.

Before defining the trustworthiness weights, we normalize these accuracy
estimates with respect to a baseline that puts them on a meaningful
scale. This baseline is based on the estimated accuracy
\(A_{\mathrm{MLC}}\) of always predicting the overall most labeled class
across all examples' annotations
\(Y_{\mathrm{MLC}} := \arg \max_k \sum_{ij} \mathbb{1}(Y_{ij} = k)\) ,
i.e.~the class selected the most by the annotators across all examples.
This accuracy is also estimated on only the subset of examples that have
more than one annotator, \(\mathcal{I}_+\) defined in (2).

\[
A _ {\mathrm{MLC}} = \frac {1}{| \mathcal {I} _ {+} |} \sum_ {i \in \mathcal {I} _ {+}} \mathbb {1} \left(Y _ {\mathrm{MLC}} = \widehat {Y} _ {i}\right) \tag {6}
\]

Adopting this most-labeled-class-accuracy as a baseline, we compute
normalized versions of our estimates for: each annotator's agreement
with other annotators and the adjusted accuracy of the model.

\[
w _ {j} = 1 - \frac {1 - s _ {j}}{1 - A _ {\mathrm{MLC}}} \tag {7}
\]

\[
w _ {\mathcal {M}} = \left(1 - \frac {1 - A _ {\mathcal {M}}}{1 - A _ {\mathrm{MLC}}}\right) \cdot \sqrt {\frac {1}{n} \sum_ {i} | \mathcal {J} _ {i} |} \tag {8}
\]

CROWDLAB uses \(w_{j}\) and \(w_{M}\) to weight our annotators and
classifier model in its weighted ensemble of predictors. Each
trustworthiness weight can thus be understood as 1 minus the (estimated)
relative error of the corresponding predictor. Such normalized-error
based weighting is commonly employed to combine predictors in model
averaging.

\section{2.2. Scoring Annotator
Quality}\label{scoring-annotator-quality}

Beyond estimating consensus labels and their quality, we consider
ranking which annotators provide the best/worst labels. Here are methods
to get an overall quality score \(a_{j} \in [0, 1]\) summarizing each
annotator's accuracy/skill.

Agreement (Monarch, 2021b). A simple score is the empirical accuracy of
each annotator's labels with respect to majority-vote consensus labels.
Examples with one annotation are not considered in this calculation to
reduce bias.

\[
a _ {j} = \frac {1}{| \mathcal {I} _ {j , +} |} \sum_ {i \in \mathcal {I} _ {j, +}} \mathbb {1} (Y _ {i j} = \widehat {Y} _ {i})
\]

where \(I_{j,+} := I_j \cap I_+ = \{i \in I_j : |J_i| > 1\}\) (9)

Label Quality Score (Kuan and Mueller, 2022). Agreement scores rate
annotators solely based on labeling statistics. We can also rely on our
classifier predictions \(\hat{p}_{M}\) to rate the average quality of
all labels provided by one annotator.

\[
a _ {j} = \frac {1}{\left| \mathcal {I} _ {j} \right|} \sum_ {i \in \mathcal {I} _ {j}} L \left(Y _ {i j}, \widehat {p} _ {\mathcal {M}} \left(Y _ {i} \mid X _ {i}\right)\right) \tag {10}
\]

CROWDLAB. Our method takes into account both the label quality score of
each annotator's labels (computed based on our classifier), as well as
the agreement between each annotator's label and the CROWDLAB consensus
label. As in (10), we estimate an average label quality score of labels
given by each annotator, but here using estimated class probabilities
\(\hat{p}_{\mathrm{CR}}\) from CROWDLAB in Sec. 2.1:

\[
Q _ {j} = \frac {1}{| \mathcal {I} _ {j} |} \sum_ {i \in \mathcal {I} _ {j}} L \left(Y _ {i j}, \hat {p} _ {\mathrm{CR}} \left(Y _ {i} \mid X _ {i}, \left\{Y _ {i j} \right\}\right)\right) \tag {11}
\]

Next, we compute each annotator's agreement with consensus among
examples with over one annotation.

\[
A _ {j} = \frac {1}{\left| \mathcal {I} _ {j , +} \right|} \sum_ {i \in \mathcal {I} _ {j, +}} \mathbb {1} \left(Y _ {i j} = \widehat {Y} _ {i}\right) \tag {12}
\]

Here \(I_{j,+}\) is defined in (9) and the consensus labels
\(\widehat{Y}_{i}\) are established via the CROWDLAB method from Sec.
2.1. Since CROWDLAB is an effective method to estimate consensus labels
\(\widehat{Y}_{i}\) , one might wonder why \(A_{j}\) alone from (12)
does not produce the best estimate of annotator quality. One reason is
that \(A_{j}\) fails to account for our confidence in each consensus
label and how individual annotators deviate from consensus. If two
annotators exhibit the same overall rate of agreement with the consensus
labels, we should favor the annotator whose deviations from consensus
are predicted to be likely classes by the classifier and tend to occur
for examples with lower consensus quality score.

Therefore we base CROWDLAB's annotator quality score on a weighted
average between \(A_{j}\) and \(Q_{j}\) . Using the model/annotator
weights \(w_{M}, w_{j}\) computed by CROWD-LAB in (7) and (8)), we find
a single aggregate weight to compare all annotators against the
classifier.

\[
\bar {w} = \frac {w _ {\mathcal {M}}}{w _ {\mathcal {M}} + w _ {0}}
\]

where
\(w_{0}=\frac{1}{nm}\sum_{i=1}^{n}\sum_{j=1}^{m}w_{j}\cdot|\mathcal{J}_{i}|\)
(13)

Here \(\bar{w}\) is shared across all annotators. It represents the
(estimated) relative trustworthiness of our classifier against the
average annotator. A quality score for each annotator is finally
computed via a weighted average of: the label quality score and the
annotator agreement with consensus labels:

\[
a _ {j} = \bar {w} Q _ {j} + (1 - \bar {w}) A _ {j} \tag {14}
\]

\section{3. Related Work}\label{related-work}

Prior work for estimating (1)-(3) from multi-annotator datasets has
fallen into two camps. The first camp relies on statistical generative
models that only account for the observed annotator statistics
(Carpenter, 2008). Like CROWDLAB, the second camp of approaches also
models feature-label relationships, but does so via linear models (Jin
et al., 2017), autoencoders (Liu et al., 2021a), or classifiers fit to
soft labels in an iterative manner (Raykar et al., 2010; Khetan et al.,
2018; Rodrigues and Pereira, 2018; Platanios et al., 2020; Liu et al.,
2021a). These methods cannot utilize an arbitrary classifier (trained
via any procedure), and they are more specialized and complex than
CROWDLAB. Due to this complexity, approaches from the former camp remain
much more popular in practical applications (Toloka).

The following sections describe existing baseline methods to estimate
consensus/annotator quality that our subsequent experiments compare
CROWDLAB against. We focus our comparison on approaches which are
either: commonly used in practice, or able to utilize any classifier to
produce better estimates (rather than approaches that are restricted to
a specific type of model or non-standard training procedure).

\section{3.1. Baseline Consensus Quality
Scores}\label{baseline-consensus-quality-scores}

Dawid-Skene (Dawid and Skene, 1979). This Bayesian method specifies a
generative model of the dataset annotations. It employs iterative
expectation-maximization (EM) to estimate each annotator's error rates
in a class-specific manner. A key estimate in this approach is
\(\widehat{p}_{\mathrm{DS}}(Y_{i} \mid \{Y_{ij}\})\) , the posterior
probability vector of the true class \(Y_{i}\) for the ith example,
given the dataset annotations \(\{Y_{ij}\}\) .

Define \(\pi_{k,\ell}^{(j)}\) as the probability that annotator \(j\)
labels an example as class \(\ell\) when the true label of that example
is \(k\) . This individual class confusion matrix for each annotator

serves as the likelihood function of the Dawid-Skene generative model.
The Dawid-Skene posterior distribution for a particular example is
computed by taking product of each annotator's likelihood and some prior
distribution \(\pi_{prior}\) .

\[
\widehat {p} _ {\mathrm{DS}} \left(Y _ {i} \mid \left\{Y _ {i j} \right\}\right) \propto \pi_ {\text { prior }} \cdot \prod_ {j \in \mathcal {J} _ {i}} \pi_ {k, Y _ {i j}} ^ {(j)} \tag {15}
\]

Our work follows conventional practice taking the prior to be the
(empirical) marginal distribution of given labels over the full dataset.
A natural consensus quality score is the label quality score of the
consensus label under the Dawid-Skene posterior class probabilities:
\(q_{i} = L(\widehat{Y}_{i}, \widehat{p}_{\mathrm{DS}}(Y_{i} \mid \{Y_{ij}\}))\)
.

GLAD (Generative model of Labels, Abilities and Difficulties) (Whitehill
et al., 2009). Specifying a more complex generative model of dataset
annotations than Dawid-Skene, this Bayesian approach also employs
iterative EM steps. GLAD additionally infers \(\alpha\) , the expertise
of each annotator and \(\beta\) , the difficulty of each example. GLAD's
likelihood is based on the following probability that an annotator
chooses the same class as the consensus label:

\[
p \left(Y _ {i j} = \widehat {Y} _ {i} \mid \alpha_ {j}, \beta_ {i}\right) = \frac {1}{1 + e ^ {- \alpha_ {j} \beta_ {i}}} \tag {16}
\]

Like Dawid-Skene, GLAD uses the data likelihood to estimate the
posterior probability of the true class \(Y_{i}\) for the ith example:
\(\widehat{p}_{\mathrm{G}}(Y_{i} \mid \{Y_{ij}\})\) . Here we use the
same standard prior as for Dawid-Skene. Again a consensus quality score
can naturally be obtained via the label quality score computed with
respect to the GLAD posterior class probabilities:
\(q_{i} = L(\widehat{Y}_{i}, \widehat{p}_{\mathrm{G}}(Y_{i} \mid \{Y_{ij}\}))\)
.

While other Bayesian annotation models exist (Kara et al., 2015; Hovy et
al., 2013; Carpenter, 2008), Dawid-Skene and GLAD are often used in
practice (Toloka; Monarch, 2021c) and perform strongly in empirical
benchmarks (Sheshadri and Lease, 2013; Paun et al., 2018; Sinha et al.,
2018).

Dawid-Skene with Model (Monarch, 2021a). Although very popular, the
Dawid-Skene and GLAD methods do not utilize a classifier at all. Thus
they struggle with sparsely labeled examples. A straightforward
adaptation of these methods to incorporate a classifier is to produce
class predictions for each example (predict hard labels rather than
probability vectors), and treat these predicted labels as if they were
the outputs from an additional annotator (Monarch, 2021a). Because
methods like Dawid-Skene and GLAD automatically adjust for estimated
annotator quality, they should theoretically account for the
classifier's strengths/weaknesses.

We augment Dawid-Skene in this way by: adding the model's predicted
labels as an additional annotator (for every example), and then
computing consensus quality scores using the same Dawid-Skene method
described above. The resulting posterior is now a function of the
example's feature values as well (since classifier predictions depend on
\(X_{i}\) ). GLAD with Model (Monarch, 2021a). We follow the same
approach to adapt GLAD to leverage the classifier: First add the model's
predicted label for each example as labels from one additional
annotator, and then compute the consensus quality score using the GLAD
method described above. Jin et al.~(2017) proposed a different extension
of GLAD to account for feature information, but their approach does not
accommodate arbitrary classifiers (eg. not suitable for images).

Empirical Bayes. While the previous two methods do not account for the
classifier's confidence in its individual predictions, we consider an
alternative adaptation of Dawid-Skene that does. This method treats the
model's prediction as a per-example prior distribution and the
annotators' labels as observations to compute
\(\widehat{p}_{\mathrm{EB}}(Y_i \mid X_i, \{Y_{ij}\})\) , the posterior
probability of the true class \(Y_i\) for the \(i\) th example, given
the dataset annotations \(\{Y_{ij}\}\) and an example-specific prior
based on the feature values \(X_i\) . The likelihood function for each
annotator is defined by the class confusion matrix estimated via the
Dawid-Skene algorithm. Using the classifier-derived prior distribution
and likelihoods, we can compute an Empirical Bayes posterior in the same
way outlined for Dawid-Skene:

\[
\widehat {p} _ {\mathrm{EB}} \left(Y _ {i} \mid X _ {i}, \left\{Y _ {i j} \right\}\right) \propto \widehat {p} _ {\mathcal {M}} \left(Y _ {i} \mid X _ {i}\right) \cdot \prod_ {j \in \mathcal {J} _ {i}} \pi_ {k, Y _ {i j}} ^ {(j)} \tag {17}
\]

and compute a consensus quality score in the same manner:
\(q_{i} = L(\hat{Y}_{i}, \hat{p}_{\mathrm{EB}}(Y_{i} \mid X_{i}, \{Y_{ij}\}))\)
.

Some have considered iterative variants of this hybrid
generative/discriminative approach, in which the classifier is retrained
to fit the resulting posterior and the above process is repeated with
the new classifier (Raykar et al., 2010; Khetan et al., 2018; Rodrigues
and Pereira, 2018; Platanios et al., 2020). This however requires a
classifier that can be iteratively trained over many rounds and also fit
to soft labels, rather than a standard classification model.

Active Label Cleaning (Bernhardt et al., 2022). Also utilizing a trained
classifier, this recently proposed method scores multi-annotator
consensus quality by subtracting the cross-entropy between classifier
predicted probabilities and individual annotations by the entropy of the
former.

\[
\begin{array}{l} q _ {i} = - \sum_ {k = 1} ^ {K} \widehat {p} _ {\text { emp }} \left(Y _ {i} = k \mid \left\{Y _ {i j} \right\} _ {j \in \mathcal {J} _ {i}}\right) \cdot \log \widehat {p} _ {\mathcal {M}, i, k} \\ - \left(- \sum_ {k = 1} ^ {K} \hat {p} _ {\mathcal {M}, i, k} \cdot \log \hat {p} _ {\mathcal {M}, i, k}\right) \tag {18} \\ \end{array}
\]

Here we abbreviate
\(\widehat{p}_{\mathcal{M},i,k} := \widehat{p}_{\mathcal{M}}(Y_i = k \mid X_i)\)
, and
\(\widehat{p}_{\mathrm{emp}}(Y_i = k \mid \{Y_{ij}\}_{j \in \mathcal{J}_i})\)
is the overall empirical distribution of class labels amongst the
annotations for a particular example. Like CROWDLAB, this approach
accounts for classifier confidence and all individual annotations. It
lacks CROWDLAB's ability to adjust for how trustworthy the individual
annotators and classifier are.

\section{3.2. Baseline Annotator Quality
Scores}\label{baseline-annotator-quality-scores}

Dawid-Skene (Dawid and Skene, 1979). We follow the conventional use of
Dawid-Skene to rate a particular annotator via the probability that they
agree with the true label. This is directly estimated for each possible
true label as part of the per-annotator class confusion matrix used by
the Dawid-Skene method (see Sec. 3.1). Thus one can score each annotator
using the trace of their confusion matrix.

\[
a _ {j} = \frac {1}{K} \sum_ {k = 1} ^ {K} \pi_ {k, k} ^ {(j)} \tag {19}
\]

GLAD (Whitehill et al., 2009). Expertise of each annotator as estimated
by GLAD method (see Sec. 3.1): \(a_{j} = \alpha_{j}\) .

Dawid-Skene with Model (Monarch, 2021a). Add the classifier's predicted
labels as an additional annotator (who labeled every example). Then
score each real annotator's quality using the Dawid-Skene method above.

GLAD with Model (Monarch, 2021a). Add the classifier's predicted labels
as an additional annotator (who labeled every example). Then score each
real annotator's quality using the GLAD method above.

\section{4. Why CROWDLAB can produce better estimates than other
methods}\label{why-crowdlab-can-produce-better-estimates-than-other-methods}

\begin{itemize}
\tightlist
\item
  In settings with few (or only one) labels for an example, the
  agreement/Dawid-Skene/GLAD scores become unreliable (Paun et al.,
  2018). CROWDLAB can utilize additional information provided by a
  classifier that may be able to generalize to this example (especially
  if other dataset examples with similar feature values have more
  trustworthy consensus labels, e.g.~if they received more
  annotations).\\
\item
  For examples that received a large number of annotations, CROWDLAB
  assigns less relative weight to the classifier predictions and its
  consensus quality score converges toward the observed annotator
  agreement. This quantity becomes more reliable when based on a large
  number of annotations (Paun et al., 2018), in which case relying on
  other sources of information becomes unnecessary. For examples where
  all annotations agree, an increase in the number of such annotations
  will typically correspond to an increased CROWDLAB consensus score.
  The Label Quality Score alone fails to exhibit this desirable
  property.\\
\item
  CROWDLAB uses weighted ensembling to combine annotations and
  classifier outputs. Countless prediction competitions have proven this
  to be among the most accurate/calibrated ways to combine different
  predictors.\\
  • Methods like Dawid-Skene estimate \(K \times K\) confusion
\end{itemize}

matrices per annotator, which may be statistically challenging when some
annotators provide few labels (Paun et al., 2018). CROWDLAB merely
estimates a single likelihood parameter P shared across all
classes/annotators in (3) as well as a single per annotator statistic
\(w_{j}\) . Both can be better estimated from a limited number of
observations.

- Generative-based methods like Dawid-Skene or GLAD are iterative
algorithms, with high computational costs when their convergence is slow
(Sinha et al., 2018; Stephens, 2000). CROWDLAB does not require
iterative updates and is deterministic (for a given classifier).

\section{5. Experiments}\label{experiments}

Datasets. To evaluate various methods, we employ real-world
multi-annotator data with naturally occurring label errors. We run three
benchmarks based on different subsets of the CIFAR-10H data (Peterson et
al., 2019) which we call: Hardest, Uniform, Complete (see Appendix A.1
for details and Appendix B and C for additional results). CIFAR-10H
contains multiple labels for images in the CIFAR-10 test set (Krizhevsky
and Hinton, 2009), obtained from a large set of new human annotators. As
a source of ground truth labels, we simply use the corresponding labels
for each image from the original CIFAR-10 dataset (Krizhevsky and
Hinton, 2009). Northcutt et al.~(2021a) found the original CIFAR-10
labels to contain few errors in verification studies, and they have been
adopted as ground truth labels in other research as well (Kuan and
Mueller, 2022).

Models. To study how methods perform across different types of
classifiers with varying accuracy, we applied every method twice, once
using a ResNet-18 classifier (He et al., 2016) and another time with a
Swin Transformer model (Liu et al., 2021b). Both classifiers are trained
on the same data (majority-vote consensus labels) in the same manner.
Here the Swin Transformer represents a high quality model, whereas
ResNet-18 represents a less accurate model (that is still commonly used
in practice).

Metrics. To measure each of our three previously stated estimation
tasks, we employ the following metrics:

\begin{enumerate}
\def\labelenumi{\arabic{enumi}.}
\tightlist
\item
  To evaluate how well methods can estimate consensus labels from
  multiply-annotated data, we measure the accuracy of the inferred
  consensus label for each example against its ground truth label.\\
\item
  To evaluate how well methods can estimate the quality of each given
  consensus label, we compare estimated quality score \(q_{i}\) for each
  example against a binary target indicating whether or not the
  consensus label matches the ground truth label. If our goal is to
\end{enumerate}

use the quality scores to flag those examples whose consensus label is
currently incorrect, this is a form of information retrieval (Kuan and
Mueller, 2022). Thus our consensus quality scores are evaluated via
precision/recall metrics: AUROC, AUPRC, and Lift at various cutoffs
(which is directly proportional to Precision@T). To focus our evaluation
purely on the estimation of label quality, throughout this section, we
use each method to estimate quality scores for a single set of consensus
labels established via majority vote. We always score the same of
consensus labels here because our above evaluation already quantifies
how good the consensus labels are from different methods, and we do not
want this to confound our evaluation of how well different methods can
estimate label quality.

\begin{enumerate}
\def\labelenumi{\arabic{enumi}.}
\setcounter{enumi}{2}
\tightlist
\item
  To evaluate how well methods can estimate the quality of each
  annotator, we measure the Spearman correlation between \(a_j\) and
  \(\mathrm{ACC}_j\) over all annotators \(j\) , where: \(a_j\) denotes
  our estimated annotator quality score (Sec. 2.2) and
  \(\mathrm{ACC}_j\) denotes the accuracy of the \(j\) -th annotator's
  chosen labels with respect to the ground truth labels (considering
  only the subset of examples labeled by annotator \(j\) ). A method
  that achieves high Spearman correlation must produce annotator quality
  scores that are lower for those annotators whose labels tend to be
  wrong the most often.
\end{enumerate}

\section{6. Results}\label{results}

Figures 2, S1, S2, and Tables 1, S2, S3 demonstrate that CROWDLAB
overall performs the best across our evaluations for consensus and
annotator quality scores, and also typically produces the most accurate
consensus labels. For most methods considered in this paper, all
evaluation metrics improve when used with the Swin Transformer
vs.~ResNet-18 model. This illustrates how a better classifier can be
utilized to get more improvement in consensus labels and
consensus/annotator quality estimates. Effective methods for
multi-annotator analysis must remain compatible with future innovations
in classifier technology.

Considering only classifier predictions and consensus labels (rather
than individual annotator information), the Label Quality Score also
effectively estimates consensus quality when we have an accurate model
(Swin Transformer). Predictions from a strong classifier suffice to
estimate label quality without additional information provided by
individual annotators (Kuan and Mueller, 2022). However Label Quality
Score performs worse than other methods with a lower accuracy classifier
(ResNet-18). This demonstrates the value of accounting for the
individual annotations and overall model accuracy in CROWDLAB, which
performs well relative to other methods regardless of the classifier's
accuracy. Treating the classifier as an additional annotator for the
Dawid-Skene and GLAD methods improves their performance, but not enough
to match CROWDLAB, which better accounts for the classifier's
confidence. While the Empirical Bayes method also accounts for
classifier confidence to augment Dawid-Skene, it similarly unable to
match CROWDLAB, demonstrating why our method considers how much to weigh
the model based on its estimated trustworthiness relative to the
annotators.

In Appendix E, we also compare CROWDLAB against a variant of our method
which lacks the per-annotator quality estimation (i.e.~all annotator
weights \(w_{j}\) are equal). Empirically, this variant underperforms as
it estimates too little information about the annotators. On another
Uniform dataset in which there are 1-5 annotations for each example
occurring with equal frequencies, CROWDLAB is able to produce better
estimates for tasks (1)-(3) than the other methods considered here
(results in Appendix B). On another Complete dataset with many more (
\(\sim 50\) ) annotations per example, such that simple annotator
agreement and majority vote produce highly accurate estimates, CROWDLAB
retains its strong performance compared to other methods (results in
Appendix C). In Appendix D, we run all methods with an unrealistically
accurate classifier on all datasets. This setting favors the Label
Quality Score, but we find that CROWDLAB still outperforms the other
methods. This breadth of settings highlights the utility of CROWDLAB
across a wide range of applications involving fair/stellar classifier
models and varying numbers of data annotations.

\section{7. Discussion}\label{discussion}

Unlike other ways to utilize classifiers with crowdsourcing algorithms,
CROWDLAB considers a model's estimated confidence and how accurate it is
relative to individual annotators. Methods such as Dawid-Skene (or GLAD)
with Model account for the model predictions, but fail adjust for model
confidence and accuracy which is properly done by CROWDLAB. Our proposed
methodology is compatible with any classifier and training strategy,
ensuring its out-of-the-box performance will improve as new models and
training tricks are invented. This is vital as classifier technology
continues to improve, and ensures CROWDLAB can be applied to diverse
data (image, text, tabular, audio, etc).

The efficacy of CROWDLAB depends on being able to train a performant
classifier, unlike generative models of annotator statistics.
Fortunately, training good classifiers is easy with modern AutoML
(Erickson et al., 2020) and techniques for calibration, data
augmentation, and transfer learning (Thulasidasan et al., 2019). As with
most classification projects, CROWDLAB users should remain wary of
overconfident model predictions with limited ability to generalize,
which may lead to overly optimistic estimates of quality. Ensuring
predictions are out-of-sample helps mitigate this.

\textit{[Image omitted: images/9844346172f843876aaead0d79802d4658c3aefe3cec22f13a9e02a16a36bc49.jpg]}

bar

AUROC for Consensus Quality Score \textbar{} Method \textbar{} AUROC
(Blue) \textbar{} AUROC (Orange) \textbar{} \textbar{} :--- \textbar{}
:--- \textbar{} :--- \textbar{} \textbar{} Agreement \textbar{} 0.68
\textbar{} 0.67 \textbar{} \textbar{} Dawid-Skene \textbar{} 0.75
\textbar{} 0.74 \textbar{} \textbar{} GLAD \textbar{} 0.88 \textbar{}
0.88 \textbar{} \textbar{} Dawid-Skene with Model \textbar{} 0.87
\textbar{} 0.89 \textbar{} \textbar{} GLAD with Model \textbar{} 0.92
\textbar{} 0.93 \textbar{} \textbar{} Empirical Bayes \textbar{} 0.91
\textbar{} 0.93 \textbar{} \textbar{} Active Label Cleaning \textbar{}
0.81 \textbar{} 0.86 \textbar{} \textbar{} Label Quality Score
\textbar{} 0.90 \textbar{} 0.94 \textbar{} \textbar{} CROWDLAB
\textbar{} 0.94 \textbar{} 0.96 \textbar{}

\textit{[Image omitted: images/198d6ee0addc074e9694203b7d106f8e0822e964442e53c26c1d9d64b9b68c33.jpg]}

bar

AUPRC for Consensus Quality Score \textbar{} Method \textbar{} AUPRC
\textbar{} \textbar{} :--- \textbar{} :--- \textbar{} \textbar{}
Agreement \textbar{} 0.15 \textbar{} \textbar{} Dawid-Skene \textbar{}
0.37 \textbar{} \textbar{} GLAD \textbar{} 0.51 \textbar{} \textbar{}
Dawid-Skene with Model \textbar{} 0.28 \textbar{} \textbar{} GLAD with
Model \textbar{} 0.40 \textbar{} \textbar{} Empirical Bayes \textbar{}
0.36 \textbar{} \textbar{} Active Label Cleaning \textbar{} 0.32
\textbar{} \textbar{} Label Quality Score \textbar{} 0.45 \textbar{}
\textbar{} CROWDLAB \textbar{} 0.54 \textbar{}

\textit{[Image omitted: images/89b34891adaad8894241433d174d1536891eb6af8817b8069873723f3e8c2c34.jpg]}

bar

Spearman Correlation for Annotator Quality \textbar{} Method \textbar{}
Spearman Correlation (Blue) \textbar{} Spearman Correlation (Orange)
\textbar{} \textbar{} :--- \textbar{} :--- \textbar{} :--- \textbar{}
\textbar{} Agreement \textbar{} 0.69 \textbar{} 0.71 \textbar{}
\textbar{} Dawid-Skene \textbar{} 0.59 \textbar{} 0.60 \textbar{}
\textbar{} GLAD \textbar{} 0.65 \textbar{} 0.64 \textbar{} \textbar{}
Dawid-Skene with Model \textbar{} 0.66 \textbar{} 0.72 \textbar{}
\textbar{} GLAD with Model \textbar{} 0.69 \textbar{} 0.72 \textbar{}
\textbar{} Label Quality Score \textbar{} 0.51 \textbar{} 0.68
\textbar{} \textbar{} CROWDLAB \textbar{} 0.72 \textbar{} 0.75
\textbar{}

\textit{[Image omitted: images/a7be3752e39936d051d638e73dae35cf379ccdf4f75c9ec5c640994cc602d98d.jpg]}

bar

Consensus Label Accuracy \textbar{} Model \textbar{} Accuracy \textbar{}
\textbar{} :--- \textbar{} :--- \textbar{} \textbar{} Majority Vote
(Agreement) \textbar{} 0.958 \textbar{} \textbar{} Dawid-Skene
\textbar{} 0.945 \textbar{} \textbar{} GLAD \textbar{} 0.945 \textbar{}
\textbar{} Dawid-Skene with Model \textbar{} 0.965 \textbar{} \textbar{}
GLAD with Model \textbar{} 0.959 \textbar{} \textbar{} CROWDLAB
\textbar{} 0.971 \textbar{}

\textit{[Image omitted: images/253f1cb93205c4746e6e60839a62e9fb3d7c730cacff2a87cfaa4987f8e316d8.jpg]}\\
Figure 2. Benchmarking methods to estimate consensus labels, their
quality, and annotator quality on the Hardest multi-annotator dataset.

Model

Method

Lift @ 10

Lift @ 50

Lift @ 100

Lift @ 300

Lift @ 500

ResNet-18

Agreement

4.87

5.84

6.33

5.27

4.72

ResNet-18

Dawid-Skene

12.89

11.79

13.26

10.74

8.51

ResNet-18

GLAD

14.6

15.69

15.88

13.93

9.96

ResNet-18

Dawid-Skene with Model

12.54

11.47

8.6

5.56

5.09

ResNet-18

GLAD with Model

14.67

13.69

14.43

11.25

10.81

ResNet-18

Empirical Bayes

12.17

12.17

11.68

11.11

10.27

ResNet-18

Active Label Cleaning

17.03

19.95

16.3

10.22

7.88

ResNet-18

Label Quality Score

19.46

21.41

19.22

13.38

10.22

ResNet-18

CROWDLAB

24.33

22.38

17.76

14.27

11.82

Swin

Agreement

2.51

6.03

6.28

4.86

4.17

Swin

Dawid-Skene

12.89

14.36

14.55

10.99

8.66

Swin

GLAD

14.6

15.69

15.88

14.11

10.04

Swin

Dawid-Skene with Model

14.16

16.43

12.46

9.16

7.99

Swin

GLAD with Model

8.7

17.39

17.1

15.36

11.71

Swin

Empirical Bayes

12.56

9.55

11.06

11.81

11.76

Swin

Active Label Cleaning

25.13

22.11

21.36

12.81

9.25

Swin

Label Quality Score

25.13

22.61

21.86

16.92

12.81

Swin

CROWDLAB

25.13

24.62

20.85

17.76

13.82

Table 1. Evaluating the precision of various consensus quality scoring
methods on the Hardest dataset. Directly proportional to Precision@T,
Lift@T reports what fraction of the top-T ranked consensus labels are
actually incorrect, normalized by the fraction of incorrect consensus
labels expected for a random set of examples.

\section{References}\label{references}

M. Bernhardt, D. C. Castro, R. Tanno, A. Schwaighofer, K. C. Tezcan, M.
Monteiro, S. Bannur, M. P. Lungren, A. Nori, B. Glocker, et al.~Active
label cleaning for improved dataset quality under resource constraints.
Nature communications, 13(1):1--11, 2022.\\
B. Carpenter. Multilevel bayesian models of categorical data annotation.
2008.\\
J. Davis and M. Goadrich. The relationship between precision-recall and
ROC curves. In International Conference on Machine learning, pages
233--240, 2006.\\
A. P. Dawid and A. M. Skene. Maximum likelihood estimation of observer
error-rates using the EM algorithm. Journal of the Royal Statistical
Society. Series C (Applied Statistics), 28(1):20--28, 1979.\\
N. Erickson, J. Mueller, A. Shirkov, H. Zhang, P. Larroy, M. Li, and A.
Smola. AutoGluon-Tabular: Robust and accurate AutoML for structured
data. arXiv preprint arXiv:2003.06505, 2020.\\
R. Fakoor, T. Kim, J. Mueller, A. J. Smola, and R. J. Tibshirani.
Flexible model aggregation for quantile regression. arXiv preprint
arXiv:2103.00083, 2021.\\
K. He, X. Zhang, S. Ren, and J. Sun. Deep residual learning for image
recognition. In Proceedings of the IEEE Conference on Computer Vision
and Pattern Recognition, 2016.\\
Hivemind and Cloudfactory. Crowd vs.~managed team: A study on quality
data processing at scale. URL
https://go.cloudfactory.com/hubfs/02-Contents/3-Reports/Crowd-vs-Managed-Team-Hivemind-Study.pdf.\\
D. Hovy, T. Berg-Kirkpatrick, A. Vaswani, and E. Hovy. Learning whom to
trust with MACE. In Proceedings of the 2013 Conference of the North
American Chapter of the Association for Computational Linguistics: Human
Language Technologies, 2013.\\
Y. Jin, M. Carman, D. Kim, and L. Xie. Leveraging side information to
improve label quality control in crowdsourcing. In Fifth AAAI Conference
on Human Computation and Crowdsourcing, 2017.\\
Y. E. Kara, G. Genc, O. Aran, and L. Akarun. Modeling annotator
behaviors for crowd labeling. Neurocomputing, 160:141--156, 2015.\\
D. Karger, S. Oh, and D. Shah. Iterative learning for reliable
crowdsourcing systems. Advances in Neural Information Processing
Systems, 2011.

A. Khetan, Z. C. Lipton, and A. Anandkumar. Learning from noisy
singly-labeled data. In International Conference on Learning
Representations, 2018.\\
A. Krizhevsky and G. Hinton. Learning multiple layers of features from
tiny images. Master's thesis, Department of Computer Science, University
of Toronto, 2009.\\
J. Kuan and J. Mueller. Model-agnostic label quality scoring to detect
real-world label errors. In ICML DataPerf Workshop, 2022.\\
Y. Liu, W. Zhang, and Y. Yu. Aggregating crowd wisdom with side
information via a clustering-based label-aware autoencoder. In
Proceedings of the Twenty-Ninth International Conference on
International Joint Conferences on Artificial Intelligence, pages
1542--1548, 2021a.\\
Z. Liu, Y. Lin, Y. Cao, H. Hu, Y. Wei, Z. Zhang, S. Lin, and B. Guo.
Swin transformer: Hierarchical vision transformer using shifted windows.
In Proceedings of the IEEE/CVF International Conference on Computer
Vision, 2021b.\\
R. Monarch. Treating model predictions as a single annotator. In
Human-in-the-loop machine learning. Manning Publications, 2021a.\\
R. Monarch. Human-in-the-loop machine learning. Manning Publications,
2021b.\\
R. Monarch. Quality control for data annotation. In Human-in-the-loop
machine learning. Manning Publications, 2021c.\\
Q. Nguyen, H. Valizadegan, and M. Hauskrecht. Learning classification
models with soft-label information. Journal of the American Medical
Informatics Association, 21(3):501--508, 2014.\\
C. G. Northcutt, A. Athalye, and J. Mueller. Pervasive label errors in
test sets destabilize machine learning benchmarks. In Proceedings of the
35th Conference on Neural Information Processing Systems Track on
Datasets and Benchmarks, December 2021a.\\
C. G. Northcutt, L. Jiang, and I. L. Chuang. Confident learning:
Estimating uncertainty in dataset labels. Journal of Artificial
Intelligence Research, 70:1373--1411, 2021b.\\
S. Paun, B. Carpenter, J. Chamberlain, D. Hovy, U. Kruschwitz, and M.
Poesio. Comparing Bayesian models of annotation. Transactions of the
Association for Computational Linguistics, 6:571--585, 2018.\\
J. C. Peterson, R. M. Battleday, T. L. Griffiths, and O. Russakovsky.
Human uncertainty makes classification more robust. In Proceedings of
the IEEE/CVF International Conference on Computer Vision, 2019.

E. A. Platanios, M. Al-Shedivat, E. Xing, and T. Mitchell. Learning from
imperfect annotations. arXiv preprint arXiv:2004.03473, 2020.\\
V. C. Raykar, S. Yu, L. H. Zhao, G. H. Valadez, C. Florin, L. Bogoni,
and L. Moy. Learning from crowds. Journal of Machine Learning Research,
11(4), 2010.\\
F. Rodrigues and F. Pereira. Deep learning from crowds. In Proceedings
of the AAAI conference on artificial intelligence, 2018.\\
A. Sheshadri and M. Lease. Square: A benchmark for research on computing
crowd consensus. In First AAAI conference on human computation and
crowdsourcing, 2013.\\
V. B. Sinha, S. Rao, and V. N. Balasubramanian. Fast Dawid-Skene: A fast
vote aggregation scheme for sentiment classification. arXiv preprint
arXiv:1803.02781, 2018.

M. Stephens. Dealing with label switching in mixture models. Journal of
the Royal Statistical Society: Series B (Statistical Methodology),
62(4):795--809, 2000.\\
S. Thulasidasan, G. Chennupati, J. A. Bilmes, T. Bhattacharya, and S.
Michalak. On mixup training: Improved calibration and predictive
uncertainty for deep neural networks. Advances in Neural Information
Processing Systems, 32, 2019.\\
Toloka. Crowd-kit: Computational quality control for crowdsourcing. URL
https://toloka.ai/en/docs/crowd-kit.\\
J. Whitehill, T.-f.~Wu, J. Bergsma, J. Movellan, and P. Ruvolo. Whose
vote should count more: Optimal integration of labels from labelers of
unknown expertise. In Advances in Neural Information Processing Systems,
2009.

\section{Appendix -- CROWDLAB: Supervised learning to infer consensus
labels and quality scores for data with multiple
annotators}\label{appendix-crowdlab-supervised-learning-to-infer-consensus-labels-and-quality-scores-for-data-with-multiple-annotators}

\section{A. Experiment Details}\label{a.-experiment-details}

Our experiments employ two of the most currently popular architectures
for image classification, which are intended to be representative of
different types of models one might use in practice. Training of the
Swin Transformer and ResNet classifiers was done as by Kuan and Mueller
(2022), using 5-fold cross-validation starting with ImageNet-pretrained
weights fine-tuned in each fold via AutoML (Erickson et al., 2020). When
establishing consensus labels via majority vote, we break ties for an
example by favoring the class which was annotated more often overall
across the dataset. We do not evaluate annotator quality scores from the
Active Label Cleaning method because rating annotators was left as
future work in the paper of Bernhardt et al.~(2022). Annotator quality
estimates from the Empirical Bayes approach match those from Dawid-Skene
and are also omitted from our plots.

Note that all metrics discussed in Section 5 are for evaluation purposes
only, and would not be computable in real applications of our
methodology due to a lack of ground truth labels. For evaluating
consensus quality scores, AUROC measures how well these scores are able
to differentiate correct and incorrect consensus labels. AUPRC accounts
for the precision/recall of the consensus quality scores in flagging an
incorrect consensus label, in a manner that is more sensitive to
proportion of errors in the majority-vote consensus label errors than
AUROC (Davis and Goadrich, 2006). The Lift at T metric measures how much
more likely we are to encounter an incorrect consensus label among the
top T ranked examples that have the worst consensus quality score.
Although our evaluation of consensus quality estimation is applied to
majority-vote consensus labels in Section 5, each method could be used
to estimate the quality of consensus labels derived via another
approach.

\section{A.1. Datasets}\label{a.1.-datasets}

The original CIFAR-10 dataset (Krizhevsky and Hinton, 2009) is fairly
easy to label (Northcutt et al., 2021a). Thus annotator agreement on the
complete CIFAR-10H data (Peterson et al., 2019) is unrealistically high
for a representative multi-annotator benchmark. The images are not only
easy to label, but there is also an uncommonly large number of
annotators ( \(\sim\) 50) per image in CIFAR-10H. Labeling budgets are
typically too small to have so many annotators review each example.
Hence our primary benchmark uses a subset of the CIFAR-10H annotations.
This subset starts with the 25 worst annotators and then incrementally
add annotators from worst to best (based on their accuracy
vs.~ground-truth labels) until each of the 10,000 examples have at least
1 annotation (resulting in a dataset with 511 annotators in total).
During this process, we restricted the selection of each new annotator
to add to the current subset to only those which labeled at least one
example also labeled by one of the annotators in the current subset. We
call this variation of CIFAR-10H the Hardest dataset benchmarked in this
paper, and believe it is more representative of real-world data labeling
applications, where the proportion of label errors tends to be far
higher than in CIFAR-10H and the number of annotators far lower
(Hivemind and Cloudfactory).

To ensure the robustness of our conclusions, we also evaluated all
methods on two other datasets: a Uniform subset of CIFAR-10H (only
considering some randomly chosen annotators such that each example has
between 1-5 annotations with an equal number of examples receiving 1
annotation, 2 annotations, etc.), and the complete CIFAR-10H dataset
(with all annotator labels, which is far more than typically collected
in most applications). Results for these other datasets are in Appendix
B and C, and are based on separate classifier models trained for each
dataset. In all cases, we only consider images from the test set of
CIFAR-10 (here treated as multiply-labeled training data), since these
are the only images labeled by many annotators in CIFAR-10H.

Labels predicted by

Accuracy (w.r.t. ground truth labels)

ResNet-18

0.879

Swin Transformer

0.940

Swin Transformer trained with true labels

0.948

Annotator (Average)

0.909

Table S1. Classification accuracy for the Hardest dataset achieved by
various predictors: ResNet-18 and Swin Transformer classifiers trained
on majority-vote consensus labels (i.e.~the models used in the benchmark
results of Figure 2), Swin Transformer trained on true labels, which
represents an unrealistically good classifier (see Appendix D), as well
as the average annotator in the dataset.

\section{B. Results for Uniform
Dataset}\label{b.-results-for-uniform-dataset}

To evaluate our methods in another setting, we construct a different
subset of CIFAR10-H and re-run our benchmark on this new dataset. In
this Uniform dataset, each example now has between 1 to 5 labels, where
the number of labels per example are uniformly distributed. This dataset
contains 421 annotators and 10,000 examples. Here the annotators are
just randomly selected from the CIFAR10-H pool, and are thus higher
quality than in the Hardest dataset. The following results demonstrate
that CROWDLAB is also the best method overall for this Uniform dataset.

\textit{[Image omitted: images/592b8550c9a8f017ce5357e5dff7cee5f45989c777b50628fed9cd72986e2dbb.jpg]}

bar

AUROC for Consensus Quality Score \textbar{} Method \textbar{} AUROC
(Blue) \textbar{} AUROC (Orange) \textbar{} \textbar{} :--- \textbar{}
:--- \textbar{} :--- \textbar{} \textbar{} Agreement \textbar{} 0.69
\textbar{} 0.67 \textbar{} \textbar{} Dawid-Skene \textbar{} 0.69
\textbar{} 0.69 \textbar{} \textbar{} GLAD \textbar{} 0.88 \textbar{}
0.87 \textbar{} \textbar{} Dawid-Skene with Model \textbar{} 0.89
\textbar{} 0.90 \textbar{} \textbar{} GLAD with Model \textbar{} 0.93
\textbar{} 0.94 \textbar{} \textbar{} Empirical Bayes \textbar{} 0.93
\textbar{} 0.93 \textbar{} \textbar{} Active Label Cleaning \textbar{}
0.81 \textbar{} 0.84 \textbar{} \textbar{} Label Quality Score
\textbar{} 0.90 \textbar{} 0.92 \textbar{} \textbar{} CROWDLAB
\textbar{} 0.95 \textbar{} 0.96 \textbar{}

\textit{[Image omitted: images/63a4ab58ea29fc0ff85606e3c1c1755d4f593569d817ef75edfde7a506d9bc43.jpg]}

bar

AUPRC for Consensus Quality Score \textbar{} Method \textbar{} AUPRC
(Blue) \textbar{} AUPRC (Orange) \textbar{} \textbar{} :--- \textbar{}
:--- \textbar{} :--- \textbar{} \textbar{} Agreement \textbar{} 0.13
\textbar{} 0.10 \textbar{} \textbar{} Dawid-Skene \textbar{} 0.29
\textbar{} 0.29 \textbar{} \textbar{} GLAD \textbar{} 0.35 \textbar{}
0.37 \textbar{} \textbar{} Dawid-Skene with Model \textbar{} 0.20
\textbar{} 0.25 \textbar{} \textbar{} GLAD with Model \textbar{} 0.30
\textbar{} 0.38 \textbar{} \textbar{} Empirical Bayes \textbar{} 0.28
\textbar{} 0.28 \textbar{} \textbar{} Active Label Cleaning \textbar{}
0.23 \textbar{} 0.33 \textbar{} \textbar{} Label Quality Score
\textbar{} 0.31 \textbar{} 0.45 \textbar{} \textbar{} CROWDLAB
\textbar{} 0.47 \textbar{} 0.55 \textbar{}

\textit{[Image omitted: images/95b193b1357be9225db34381a188b8f404a016a54e551e04836ea1b26b526bb9.jpg]}

bar

Spearman Correlation for Annotator Quality \textbar{} Annotator Quality
\textbar{} Spearman Correlation (Blue) \textbar{} Spearman Correlation
(Orange) \textbar{} \textbar{} :--- \textbar{} :--- \textbar{} :---
\textbar{} \textbar{} Agreement \textbar{} 0.82 \textbar{} 0.83
\textbar{} \textbar{} Dawid-Skene \textbar{} 0.76 \textbar{} 0.76
\textbar{} \textbar{} GLAD \textbar{} 0.76 \textbar{} 0.76 \textbar{}
\textbar{} Dawid-Skene with Model \textbar{} 0.75 \textbar{} 0.79
\textbar{} \textbar{} GLAD with Model \textbar{} 0.79 \textbar{} 0.81
\textbar{} \textbar{} Label Quality Score \textbar{} 0.47 \textbar{}
0.63 \textbar{} \textbar{} CROWDLAB \textbar{} 0.82 \textbar{} 0.84
\textbar{}

\textit{[Image omitted: images/c8586c294b09406da92b782c913d428fe9f9244b3451b39f0cea0f70ab4108e3.jpg]}

bar

Consensus Label Accuracy \textbar{} Model \textbar{} Accuracy \textbar{}
\textbar{} :--- \textbar{} :--- \textbar{} \textbar{} Majority Vote
(Agreement) \textbar{} 0.976 \textbar{} \textbar{} Dawid-Skene
\textbar{} 0.971 \textbar{} \textbar{} GLAD \textbar{} 0.971 \textbar{}
\textbar{} Dawid-Skene with Model \textbar{} 0.961 \textbar{} \textbar{}
GLAD with Model \textbar{} 0.978 \textbar{} \textbar{} CROWDLAB
\textbar{} 0.981 \textbar{}

ResNet-18 Swin Transformer

Figure S1. Benchmarking methods to estimate consensus labels, their
quality, and annotator quality on the Uniform dataset.

Model

Quality Method

Lift @ 10

Lift @ 50

Lift @ 100

Lift @ 300

Lift @ 500

ResNet-18

Agreement

21.19

11.86

9.75

9.04

7.12

ResNet-18

Dawid-Skene

31.69

24.65

19.01

12.09

8.52

ResNet-18

GLAD

31.69

22.54

25.7

13.03

9.23

ResNet-18

Dawid-Skene with Model

10.34

11.37

7.24

5.43

4.86

ResNet-18

GLAD with Model

17.54

18.42

15.79

13.01

12.72

ResNet-18

Empirical Bayes

12.71

20.34

18.22

13.98

11.44

ResNet-18

Active Label Cleaning

33.9

24.58

16.1

10.03

7.88

ResNet-18

Label Quality Score

33.9

26.27

22.46

12.43

9.75

ResNet-18

CROWDLAB

42.37

33.05

27.97

16.95

13.81

Swin

Agreement

8.89

8.0

7.56

7.85

6.49

Swin

Dawid-Skene

28.17

27.46

20.07

11.74

8.45

Swin

GLAD

35.21

25.35

27.11

13.03

9.23

Swin

Dawid-Skene with Model

33.33

16.3

12.22

8.89

8.15

Swin

GLAD with Model

23.47

23.47

23.0

19.87

14.74

Swin

Empirical Bayes

13.33

17.78

17.78

14.96

12.44

Swin

Active Label Cleaning

44.44

30.22

24.0

14.07

10.13

Swin

Label Quality Score

35.56

32.89

29.33

18.67

13.6

Swin

CROWDLAB

40.0

35.56

32.0

21.19

15.2

Table S2. Evaluating the precision of various consensus quality scoring
methods on the Uniform dataset. Lift@T is directly proportional to
Precision@T, and reports what fraction of the top-T ranked consensus
labels are actually incorrect normalized by the fraction of incorrect
consensus labels expected for a random set of examples.

\section{C. Results for Complete
Dataset}\label{c.-results-for-complete-dataset}

We also evaluate our methods on the full original CIFAR-10H dataset
(Peterson et al., 2019). This Complete dataset contains 2571 annotators
where each annotator labels 200 examples, such that each of the 10,000
images has approximately 50 annotations. The Complete dataset has by far
the highest number of annotations per example out of all the datasets
considered in this paper. Far less annotations per example are available
in most real-world multi-annotator datasets due to limited labeling
budgets.

With so many annotations per example, basic annotator agreement methods
are highly effective. CROWDLAB works similarly well, highlighting its
adaptive nature across different datasets with few vs.~many annotations
per example. Even in this Complete dataset where majority-vote consensus
labels should be extremely accurate, CROWDLAB consensus labels are even
more accurate (regardless whether a Swin Transformer or ResNet-18 model
is used to augment the annotators).

Even though CROWDLAB is straightforward and only estimates one parameter
per annotator ( \(w_j\) ) and two others in total ( \(P\) and
\(w_{\mathcal{M}}\) ), it performs well on this annotation-rich dataset.
This indicates CROWDLAB is not too simple, given that methods which
estimate many more parameters like Dawid-Skene and GLAD do not perform
better on this Complete dataset, even though there is no shortage of
annotations to learn from. CROWDLAB's fewer number of parameters entail
a key advantage for most datasets which have fewer annotations, and do
not cause it to lag behind richer generative methods like
Dawid-Skene/GLAD in this setting with unusually many annotations.

\textit{[Image omitted: images/1264dc1cf4d2a56d68997ce3a070789d3eb4d597dd11f3ab67775194530a479d.jpg]}

bar

AUROC for Consensus Quality Score \textbar{} Method \textbar{} Blue Bar
\textbar{} Orange Bar \textbar{} \textbar{} :--- \textbar{} :---
\textbar{} :--- \textbar{} \textbar{} Agreement \textbar{} 0.96
\textbar{} 0.96 \textbar{} \textbar{} Dawid-Skene \textbar{} 0.96
\textbar{} 0.96 \textbar{} \textbar{} GLAD \textbar{} 0.76 \textbar{}
0.76 \textbar{} \textbar{} Dawid-Skene with Model \textbar{} 0.96
\textbar{} 0.96 \textbar{} \textbar{} GLAD with Model \textbar{} 0.76
\textbar{} 0.76 \textbar{} \textbar{} Empirical Bayes \textbar{} 0.96
\textbar{} 0.96 \textbar{} \textbar{} Active Label Cleaning \textbar{}
0.79 \textbar{} 0.73 \textbar{} \textbar{} Label Quality Score
\textbar{} 0.88 \textbar{} 0.94 \textbar{} \textbar{} CROWDLAB
\textbar{} 0.96 \textbar{} 0.96 \textbar{}

\textit{[Image omitted: images/9d58318f7b30f7582383ab5bbe881f30efd6408a7cc64c8cabbbca28865c63e5.jpg]}

bar

AUPRC for Consensus Quality Score \textbar{} Method \textbar{} AUPRC
(Blue) \textbar{} AUPRC (Orange) \textbar{} \textbar{} :--- \textbar{}
:--- \textbar{} :--- \textbar{} \textbar{} Agreement \textbar{} 0.26
\textbar{} 0.25 \textbar{} \textbar{} Dawid-Skene \textbar{} 0.22
\textbar{} 0.23 \textbar{} \textbar{} GLAD \textbar{} 0.35 \textbar{}
0.36 \textbar{} \textbar{} Dawid-Skene with Model \textbar{} 0.21
\textbar{} 0.20 \textbar{} \textbar{} GLAD with Model \textbar{} 0.29
\textbar{} 0.25 \textbar{} \textbar{} Empirical Bayes \textbar{} 0.34
\textbar{} 0.32 \textbar{} \textbar{} Active Label Cleaning \textbar{}
0.04 \textbar{} 0.03 \textbar{} \textbar{} Label Quality Score
\textbar{} 0.10 \textbar{} 0.11 \textbar{} \textbar{} CROWDLAB
\textbar{} 0.31 \textbar{} 0.37 \textbar{}

\textit{[Image omitted: images/6779d30565cdb0477bfb4ae6c8c493e6d1ea2275fb994102d27d945637f3cf16.jpg]}

bar

{\def\LTcaptype{none} % do not increment counter
\begin{longtable}[]{@{}
  >{\raggedright\arraybackslash}p{(\linewidth - 4\tabcolsep) * \real{0.2329}}
  >{\raggedright\arraybackslash}p{(\linewidth - 4\tabcolsep) * \real{0.3699}}
  >{\raggedright\arraybackslash}p{(\linewidth - 4\tabcolsep) * \real{0.3973}}@{}}
\toprule\noalign{}
\begin{minipage}[b]{\linewidth}\raggedright
Annotator Quality
\end{minipage} & \begin{minipage}[b]{\linewidth}\raggedright
Spearman Correlation (Blue)
\end{minipage} & \begin{minipage}[b]{\linewidth}\raggedright
Spearman Correlation (Orange)
\end{minipage} \\
\midrule\noalign{}
\endhead
\bottomrule\noalign{}
\endlastfoot
Agreement & 0.95 & 0.95 \\
Dawid-Skene & 0.96 & 0.96 \\
GLAD & 0.88 & 0.88 \\
Dawid-Skene with Model & 0.96 & 0.96 \\
GLAD with Model & 0.88 & 0.88 \\
Label Quality Score & 0.58 & 0.72 \\
CROWDLAB & 0.96 & 0.96 \\
\end{longtable}
}

\textit{[Image omitted: images/766bb417c17aa4df7fa0a2ea74ddc2de2f3faf92027718ab571f057019733277.jpg]}

bar

Consensus Label Accuracy \textbar{} Method \textbar{} Accuracy
\textbar{} \textbar{} :--- \textbar{} :--- \textbar{} \textbar{}
Majority Vote (Agreement) \textbar{} 0.992 \textbar{} \textbar{}
Dawid-Skene \textbar{} 0.993 \textbar{} \textbar{} GLAD \textbar{} 0.992
\textbar{} \textbar{} Dawid-Skene with Model \textbar{} 0.993 \textbar{}
\textbar{} GLAD with Model \textbar{} 0.993 \textbar{} \textbar{}
CROWDLAB \textbar{} 0.993 \textbar{}

\textit{[Image omitted: images/a8be9b0c69b5cfbb2682013117c042b196965c6a78061c2f65afd5ef3a85d3e5.jpg]}\\
Figure S2. Benchmarking methods to estimate consensus labels, their
quality, and annotator quality on the Complete dataset.

Model

Quality Method

Lift @ 10

Lift @ 50

Lift @ 100

Lift @ 300

Lift @ 500

ResNet-18

Agreement

25.97

33.77

44.16

26.41

17.92

ResNet-18

Dawid-Skene

27.4

41.1

39.73

25.57

17.81

ResNet-18

GLAD

89.74

66.67

47.44

18.38

11.03

ResNet-18

Dawid-Skene with Model

28.17

36.62

38.03

26.29

17.75

ResNet-18

GLAD with Model

53.33

61.33

46.67

17.78

10.67

ResNet-18

Empirical Bayes

77.92

49.35

44.16

26.84

18.18

ResNet-18

Active Label Cleaning

0.0

10.39

12.99

8.66

8.83

ResNet-18

Label Quality Score

38.96

20.78

19.48

13.42

9.09

ResNet-18

CROWDLAB

38.96

49.35

44.16

27.71

18.44

Swin

Agreement

26.32

34.21

43.42

26.32

17.89

Swin

Dawid-Skene

41.1

41.1

39.73

25.57

17.81

Swin

GLAD

89.74

66.67

47.44

18.38

11.03

Swin

Dawid-Skene with Model

14.93

38.81

37.31

25.87

17.61

Swin

GLAD with Model

28.17

53.52

46.48

17.37

10.42

Swin

Empirical Bayes

78.95

50.0

42.11

26.32

17.89

Swin

Active Label Cleaning

0.0

0.0

5.26

7.46

6.84

Swin

Label Quality Score

26.32

21.05

21.05

17.11

13.68

Swin

CROWDLAB

65.79

65.79

48.68

28.07

18.68

Table S3. Evaluating the precision of various consensus quality scoring
methods on the Complete dataset. Lift@T is directly proportional to
Precision@T, and reports what fraction of the top-T ranked consensus
labels are actually incorrect normalized by the fraction of incorrect
consensus labels expected for a random set of examples.

\section{D. Model Trained on True CIFAR-10
Labels}\label{d.-model-trained-on-true-cifar-10-labels}

In this section, we investigate how the methods perform when utilizing a
highly accurate model. We obtain such a model by training a Swin
Transformer on the ground truth labels rather than consensus labels
estimated from the given annotations (this would not be possible in real
applications). All benchmark results presented in this section are with
respect to this unrealistically good classifier.

\textit{[Image omitted: images/610bf3cbef530b932ecaba96903ed33e67aa6a8e345a05ba912545eefba3df9f.jpg]}

bar

AUROC for Consensus Quality Score \textbar{} Method \textbar{} AUROC
(Green) \textbar{} AUROC (Orange) \textbar{} AUROC (Blue) \textbar{}
\textbar{} :--- \textbar{} :--- \textbar{} :--- \textbar{} :---
\textbar{} \textbar{} Agreement \textbar{} 0.96 \textbar{} 0.68
\textbar{} 0.65 \textbar{} \textbar{} Dawid-Skene \textbar{} 0.96
\textbar{} 0.70 \textbar{} 0.75 \textbar{} \textbar{} GLAD \textbar{}
0.96 \textbar{} 0.88 \textbar{} 0.88 \textbar{} \textbar{} Dawid-Skene
with Model \textbar{} 0.96 \textbar{} 0.91 \textbar{} 0.91 \textbar{}
\textbar{} GLAD with Model \textbar{} 0.96 \textbar{} 0.75 \textbar{}
0.93 \textbar{} \textbar{} Empirical Bayes \textbar{} 0.96 \textbar{}
0.94 \textbar{} 0.94 \textbar{} \textbar{} Active Label Cleaning
\textbar{} 0.94 \textbar{} 0.87 \textbar{} 0.89 \textbar{} \textbar{}
Label Quality Score \textbar{} 0.96 \textbar{} 0.94 \textbar{} 0.94
\textbar{} \textbar{} CROWDLAB \textbar{} 0.96 \textbar{} 0.96
\textbar{} 0.96 \textbar{}

\textit{[Image omitted: images/bb52f2f797027ec16084a8224fc1af58262835e03050996f8ad82bb038fa53e4.jpg]}

bar

AUPRC for Consensus Quality Score \textbar{} Method \textbar{} AUPRC
(Blue) \textbar{} AUPRC (Orange) \textbar{} AUPRC (Green) \textbar{}
\textbar{} :--- \textbar{} :--- \textbar{} :--- \textbar{} :---
\textbar{} \textbar{} Agreement \textbar{} 0.1 \textbar{} 0.1 \textbar{}
0.28 \textbar{} \textbar{} Dawid-Skene \textbar{} 0.42 \textbar{} 0.29
\textbar{} 0.23 \textbar{} \textbar{} GLAD \textbar{} 0.53 \textbar{}
0.35 \textbar{} 0.35 \textbar{} \textbar{} Dawid-Skene with Model
\textbar{} 0.32 \textbar{} 0.28 \textbar{} 0.23 \textbar{} \textbar{}
GLAD with Model \textbar{} 0.44 \textbar{} 0.40 \textbar{} 0.29
\textbar{} \textbar{} Empirical Bayes \textbar{} 0.40 \textbar{} 0.35
\textbar{} 0.36 \textbar{} \textbar{} Active Label Cleaning \textbar{}
0.56 \textbar{} 0.44 \textbar{} 0.1 \textbar{} \textbar{} Label Quality
Score \textbar{} 0.70 \textbar{} 0.57 \textbar{} 0.22 \textbar{}
\textbar{} CROWDLAB \textbar{} 0.70 \textbar{} 0.58 \textbar{} 0.37
\textbar{}

\textit{[Image omitted: images/c18e96e2dc22ae9e93a6a466b457fa457dd100838c03c4df5b495cca065f5cea.jpg]}

bar

{\def\LTcaptype{none} % do not increment counter
\begin{longtable}[]{@{}llll@{}}
\toprule\noalign{}
Annotator Quality & Blue Bar & Orange Bar & Green Bar \\
\midrule\noalign{}
\endhead
\bottomrule\noalign{}
\endlastfoot
Agreement & 0.72 & 0.83 & 0.95 \\
Dawid-Skene & 0.60 & 0.75 & 0.96 \\
GLAD & 0.65 & 0.75 & 0.87 \\
Dawid-Skene with Model & 0.72 & 0.78 & 0.96 \\
GLAD with Model & 0.72 & 0.80 & 0.87 \\
Label Quality Score & 0.67 & 0.63 & 0.74 \\
CROWDLAB & 0.77 & 0.84 & 0.95 \\
\end{longtable}
}

\textit{[Image omitted: images/4c817018b7200f53358979bbb4516d97b2de8d6cf0a0dccd1d9a6da50d5959ec.jpg]}

bar

Consensus Label Accuracy \textbar{} Method \textbar{} Blue Bar
\textbar{} Orange Bar \textbar{} Green Bar \textbar{} \textbar{} :---
\textbar{} :--- \textbar{} :--- \textbar{} :--- \textbar{} \textbar{}
Majority Vote (Agreement) \textbar{} 0.96 \textbar{} 0.975 \textbar{}
0.99 \textbar{} \textbar{} Dawid-Skene \textbar{} 0.945 \textbar{} 0.97
\textbar{} 0.99 \textbar{} \textbar{} GLAD \textbar{} 0.945 \textbar{}
0.97 \textbar{} 0.99 \textbar{} \textbar{} Dawid-Skene with Model
\textbar{} 0.965 \textbar{} 0.97 \textbar{} 0.99 \textbar{} \textbar{}
GLAD with Model \textbar{} 0.965 \textbar{} 0.975 \textbar{} 0.99
\textbar{} \textbar{} CROWDLAB \textbar{} 0.97 \textbar{} 0.98
\textbar{} 0.99 \textbar{}

Hardest Uniform Complete

Figure S3. Benchmarking multi-annotator methods that utilize an
unrealistically good classifier fit to true labels for each dataset.

Quality Method

Lift @ 10

Lift @ 50

Lift @ 100

Lift @ 300

Lift @ 500

Agreement

2.65

3.18

3.98

4.16

3.77

Dawid-Skene

14.73

15.47

15.1

11.48

8.95

GLAD

14.6

15.69

16.24

14.6

10.07

Dawid-Skene with Model

24.02

18.62

15.02

10.11

9.13

GLAD with Model

19.11

21.66

20.06

16.45

11.97

Empirical Bayes

5.31

7.43

10.34

12.73

12.84

Active Label Cleaning

23.87

25.46

25.2

16.18

11.03

Label Quality Score

23.87

24.93

24.93

20.07

14.64

CROWDLAB

26.53

25.99

24.14

19.19

14.8

Table S4. Evaluating the lift (i.e.~precision) of various consensus
quality scoring methods on the Hardest dataset, here employing our
unrealistically good classifier trained with true labels.

Quality Method

Lift @ 10

Lift @ 50

Lift @ 100

Lift @ 300

Lift @ 500

Agreement

8.93

10.71

8.48

7.44

6.34

Dawid-Skene

28.17

25.35

20.42

12.21

8.73

GLAD

31.69

24.65

27.82

13.03

9.23

Dawid-Skene with Model

31.01

19.38

14.73

9.95

9.69

GLAD with Model

13.95

25.12

23.72

20.78

15.44

Empirical Bayes

13.39

19.64

20.98

19.05

14.38

Active Label Cleaning

40.18

39.29

32.14

17.11

11.34

Label Quality Score

40.18

40.18

36.16

19.79

14.38

CROWDLAB

44.64

33.04

34.38

22.32

15.98

Table S5. Evaluating the lift (i.e.~precision) of various consensus
quality scoring methods on the Uniform dataset, here employing our
unrealistically good classifier trained with true labels.

Quality Method

Lift @ 10

Lift @ 50

Lift @ 100

Lift @ 300

Lift @ 500

Agreement

25.64

35.9

43.59

26.5

17.95

Dawid-Skene

27.4

38.36

39.73

25.57

17.81

GLAD

89.74

66.67

47.44

18.38

11.28

Dawid-Skene with Model

42.86

40.0

38.57

26.67

17.71

GLAD with Model

40.54

54.05

45.95

17.12

10.54

Empirical Bayes

89.74

48.72

44.87

26.92

18.21

Active Label Cleaning

38.46

23.08

19.23

14.53

11.54

Label Quality Score

64.1

51.28

38.46

18.38

12.56

CROWDLAB

89.74

61.54

42.31

26.92

18.46

Table S6. Evaluating the lift (i.e.~precision) of various consensus
quality scoring methods on the Complete dataset, here employing our
unrealistically good classifier trained with true labels.

\section{E. Variant of our Method Without Per Annotator
Weights}\label{e.-variant-of-our-method-without-per-annotator-weights}

Here we present results for a simpler variant of CROWDLAB that we also
explored, henceforth called No Perannotator Weights. The two approaches
are overall the same, except while CROWDLAB considers each annotator
individually and assigns them a separate weight \(w_{j}\) , No
Perannotator Weights aggregates all the annotators and treats them as
one ``average annotator'' to be weighed against the classifier model.
Details of the No Perannotator Weights approach are presented below.

\section{E.1. Consensus Quality
Method}\label{e.1.-consensus-quality-method}

Just as in CROWDLAB, we estimate the quality of consensus labels via the
label quality score based on estimated class probabilities. In the No
Perannotator Weights variant, these probabilities are computed via a
slightly different weighted average:

\[
\hat {p} _ {\mathrm{NPW}} \left(Y _ {i} \mid X _ {i}, \left\{Y _ {i j} \right\}\right) = \frac {w _ {\mathcal {M}} \cdot \hat {p} _ {\mathcal {M}} \left(Y _ {i} \mid X _ {i}\right) + w _ {\mathcal {A}} \cdot \hat {p} _ {\mathcal {A}} \left(Y _ {i} \mid \left\{Y _ {i j} \right\}\right)}{w _ {\mathcal {M}} + w _ {\mathcal {A}}} \tag {20}
\]

where \(w_{M} = w \cdot \frac{1}{n} \sum_{i} \sqrt{|\mathcal{J}_{i}|}\)
, \(w_{A} = (1 - w) \cdot \sqrt{|\mathcal{J}_{i}|}\) are one weight for
the model and one weight applied to all annotators. Both depend on w,
whose definition follows a similar strategy as used in CROWDLAB for
individual annotators, but here applied to their aggregate output.

First let's recall these quantities from Sec. 2.1: \(s_j\) represents
annotator \(j\) 's agreement with other annotators who labeled the same
examples and is defined in (4), \(A_j\) represents the accuracy of each
annotator's labels with respect to the majority-vote consensus label for
examples with more than one annotation and is defined in (12). In this
variant, we compute an average annotator accuracy \(\bar{A}\) by taking
the average of each annotator's accuracy weighted simply by the number
of examples each annotator labeled (rather than their estimated
trustworthiness).

\[
\bar {A} = \frac {\sum_ {j} A _ {j} \cdot | \mathcal {I} _ {j} |}{\sum_ {j} | \mathcal {I} _ {j} |}
\]

Let \(A_{\mathcal{M}}\) represent the accuracy of the model with respect
to the majority-vote consensus labels among examples with more than one
annotation, as defined in (5). We then choose our weight
\(w = A_{\mathcal{M}} / (A_{\mathcal{M}} + \bar{A})\) to balance model
accuracy vs.~that of the average annotator.

While CROWDLAB uses a separate class likelihood vector for each
annotator, this variant only considers their aggregate class likelihood

\[
\widehat {p} _ {\mathcal {A}} (Y _ {i} = k \mid \{Y _ {i j} \}) = \frac {1}{| \mathcal {J} _ {i} |} \sum_ {j \in \mathcal {J} _ {i}} P _ {j} \quad \text {where} P _ {j} = \left\{ \begin{array}{l l} s _ {j} & \text {when} Y _ {i j} = k \\ \frac {1 - s _ {j}}{K - 1} & \text {when} Y _ {i j} \neq k \end{array} \right.
\]

\section{E.2. Annotator Quality
Method}\label{e.2.-annotator-quality-method}

In the No Perannotator Weights variant, we score the quality of each
annotator via:

\[
a _ {j} = w \cdot Q _ {j} + (1 - w) \cdot A _ {j}
\]

Here \(Q_{j}\) the average label quality score of labels given by each
annotator, computed via (11) as in CROWDLAB, but here based on class
probabilities \(\widehat{p}_{NPW}\) estimated using No Perannotator
Weights defined in (20) in place of \(\widehat{p}_{CR}\) . \(A_{j}\) and
w are defined as above in Sec. E.1.

\section{E.3. Benchmarking CROWDLAB with/without per annotator
weights}\label{e.3.-benchmarking-crowdlab-withwithout-per-annotator-weights}

Ignoring the strengths and weakness of each individual annotator when
aggregating them is overall detrimental to CROWDLAB. However the
performance reduction due to this modification is surprisingly small,
given how important accounting for annotators' relative quality is
stated to be in the crowdsourcing literature (Hovy et al., 2013; Karger
et al., 2011; Kara et al., 2015; Dawid and Skene, 1979; Whitehill et
al., 2009). Rather the key aspects behind the success of CROWDLAB are
its careful consideration of: how much to trust the classifier model
vs.~the aggregate annotations along with how many annotations were
provided for each example. Studying additional variants of CROWDLAB with
either of these two pieces removed produced very poor results in our
benchmarks.

\textit{[Image omitted: images/2f454e21e6aff302cc71bcf708a90f85f4aa61cf58e55a7e5dcaa41da8c56da8.jpg]}\\
Figure S4. Benchmarking CROWDLAB with/without per annotator weights on
the Hardest dataset.

Model

Quality Method

Lift @ 10

Lift @ 50

Lift @ 100

Lift @ 300

Lift @ 500

ResNet-18

No Perannotator Weights

24.33

21.9

18.25

13.95

11.92

ResNet-18

CROWDLAB

24.33

22.38

17.76

14.27

11.82

Swin

No Perannotator Weights

25.13

23.62

20.85

17.84

14.02

Swin

CROWDLAB

25.13

24.62

20.85

17.76

13.82

Table S7. Evaluating the precision of CROWDLAB consensus quality scores
with/without per annotator weights on the Hardest dataset. Lift@T is
directly proportional to Precision@T, and reports what fraction of the
top-T ranked consensus labels are actually incorrect, normalized by the
fraction of incorrect consensus labels expected for a random set of
examples.

\end{document}
